
\subsection{能量估计}
\label{subsec:5-3-5}
\liftIndex{I196}

本小节证明命题 5.3.20。
为此，先回顾霍费尔范数和
归一化哈密顿函数的基本事实。
给定哈密顿同痕
$\{\psi_\rho\}_\rho$，
其生成哈密顿函数
$K\colon[0,1]\times M\to\mathbb R$
仅在加上一个只依赖
$\rho$ 的函数
$c\colon[0,1]\to\mathbb R$
的意义下唯一。
由于此处的 $M$ 是闭流形，
为消除这一歧义，
我们采用均值归一化：
要求光滑函数 $K$ 满足
\[
\int_M K(\rho,x)\,d\mu(x)=0
\qquad(0\le\rho\le1),
\]
其中 $d\mu$ 是 $(M,\omega)$ 的
刘维尔测度。
称这样的含时函数 $K$
为归一化哈密顿函数。
以此定义哈密顿同痕的
霍费尔长度\liftIndex{I197}\liftIndex{I198}：
\begin{equation}
\bigl\|\{\psi_\rho\}_\rho\bigr\|
:=\|K\|
=\int_0^1
\bigl(\max_x K_\rho-\min_x K_\rho\bigr)
\,d\rho.
\tag{5.3.40}\label{eq:5-3-40}
\end{equation}

事实上，霍费尔长度的定义
本身不依赖所用的归一化；
不过，下面必须注意
归一化哈密顿函数的一项性质。
这里依照 \textbf{[Oh05I]}
的论述。
令 $\psi_{\rho,t}$ 为任意双参数
哈密顿微分同胚族，并定义
哈密顿向量场 $X_{\rho,t}$、
$Y_{\rho,t}$：
\[
\frac{\partial\psi_{\rho,t}}{\partial\rho}
\circ(\psi_{\rho,t})^{-1}=X_{\rho,t},
\qquad
\frac{\partial\psi_{\rho,t}}{\partial t}
\circ(\psi_{\rho,t})^{-1}=Y_{\rho,t}.
\]
以下恒等式是关键；
其推导见 \textbf{[Ban78]}：
\begin{equation}
\frac{\partial X_{\rho,t}}{\partial t}
=\frac{\partial Y_{\rho,t}}{\partial\rho}
+[X_{\rho,t},Y_{\rho,t}].
\tag{5.3.41}\label{eq:5-3-41}
\end{equation}
以归一化哈密顿函数
$K(\rho,t,\cdot)$、
$F(\rho,t,\cdot)$
分别表示 $X_{\rho,t}$、
$Y_{\rho,t}$。\translatornote{原书实体页 331
在此处把 $F,K$ 分别配给
$\rho,t$ 方向，随后却在
(5.3.52)--(5.3.57)
把 $K,F$ 分别用作
$\rho,t$ 方向的生成函数；
原式 (5.3.42)--(5.3.43)
也不能使后文积分项相消。
这里依后文计算统一对应关系，
并据 (5.3.41) 改写两式；
见疑误 A-0305。}
于是由 (5.3.41) 可知，
存在只依赖 $(\rho,t)$ 的
$c\colon[0,1]^2\to\mathbb R$，
使得
\begin{equation}
\frac{\partial F}{\partial\rho}
(\rho,t,x)
=\frac{\partial K}{\partial t}
(\rho,t,x)
-\{F,K\}(\rho,t,x)+c(\rho,t).
\tag{5.3.42}\label{eq:5-3-42}
\end{equation}
其中 $\{F,K\}$ 是泊松括号，
即
$\{F,K\}=dF(X_K)$；
它满足
\[
[X_F,X_K]=-X_{\{F,K\}}.
\]
由刘维尔定理，
泊松括号 $\{F,K\}$
自动满足上述均值归一化。
对 (5.3.42) 在 $M$ 上积分
得 $c\equiv0$，故
\begin{equation}
\frac{\partial F}{\partial\rho}
=\frac{\partial K}{\partial t}
-\{F,K\}.
\tag{5.3.43}\label{eq:5-3-43}
\end{equation}
上述讨论归结为以下引理。

\begin{lemma}{5.3.44}\label{lem:5-3-44}
设 $\{\psi_{\rho,t}\}_{\rho,t}$
是定义在 $(\rho,t)\in[0,1]^2$
上的光滑双参数哈密顿微分同胚族，
$F,K$ 为上述相应的
归一化哈密顿函数。
则 (5.3.43) 成立。
\end{lemma}

现在把命题 5.3.20
重述为下列命题 5.3.45。
作用量泛函 $\mathcal A$
的定义见第 2.2.2 小节开头。
\translatornote{原书实体页 331
在命题 5.3.45 的输出道路空间中，
第二侧写成整个同痕
$\psi^{(1)}(L')$，
道路下标也与命题 5.3.20 不同。
这里依固定的时间 $1$ 端点
及本节既定的 $\ell_0^{\vec\psi}$
记号统一；见疑误 A-0307。}

\begin{proposition}{5.3.45}\label{prop:5-3-45}
设
$\psi^{(0)}
=\{\psi_\rho^{(0)}\}_\rho$、
$\psi^{(1)}
=\{\psi_\rho^{(1)}\}_\rho$
是端点
$\psi_1^{(0)},\psi_1^{(1)}$
分别固定的任意哈密顿同痕。
取
$[\ell_p,w]\in
\widetilde\Omega_{\ell_0}(L,L')$
及
$[\ell_q,w']\in
\widetilde\Omega_{\ell_0^{\vec\psi}}
\bigl(
\psi_1^{(0)}(L),
\psi_1^{(1)}(L')
\bigr)$，
并假设模空间 (5.3.19) 非空。
则
\begin{equation}
0\le
\mathcal A_{\ell_0^{\vec\psi}}
([\ell_q,w'])
-\mathcal A_{\ell_0}([\ell_p,w])
+\mathfrak E.
\tag{5.3.46.1}\label{eq:5-3-46-1}
\end{equation}
特别地，
\begin{equation}
\mathcal A_{\ell_0^{\vec\psi}}
([\ell_q,w'])
-\mathcal A_{\ell_0}([\ell_p,w])
\ge-\mathfrak E
=-\mathfrak E(\psi^{(0)},\psi^{(1)}).
\tag{5.3.46.2}\label{eq:5-3-46-2}
\end{equation}
事实上，可以取
\begin{equation}
\mathfrak E
=\bigl\|\{\psi_\rho^{(0)}\}_\rho\bigr\|
+\bigl\|\{\psi_\rho^{(1)}\}_\rho\bigr\|,
\tag{5.3.47}\label{eq:5-3-47}
\end{equation}
其中两项是 (5.3.40)
定义的霍费尔长度。
\end{proposition}

\begin{proof}
取非空模空间 (5.3.19) 的一个元素
$(\Sigma;u;
\vec\tau^{(0)},\vec\tau^{(1)})$。
由于 $u$ 满足以
$\omega$ 相容的殆复结构族
$\{J_{\tau,t}\}_{\tau,t}$
为参数的方程 (5.3.18.1)，
有 $\int u^*\omega\ge0$。
证明中只有这一处用到
$u$ 是柯西--黎曼方程的解；
用物理学的说法，
$\int u^*\omega\ge0$
是在壳上成立的。
以下计算则可在离壳情形进行。

现令 $u\colon\mathbb R\times[0,1]\to M$
为任意有限能量的光滑映射，
满足给定的拉格朗日边界条件，
并具有第 4.1 节末所述
足以紧化的无穷远衰减性质。
这就是说，复合映射
\[
\bar u(\rho,t)
=u(\chi^{-1}(\rho),t)
\colon(0,1)\times[0,1]\to M
\]
可连续延拓至 $[0,1]^2$。
例如，对干净相交的
$L,L'$，方程 (5.3.18.1)
的任何有限能量解
都具有此衰减性质。

先回顾命题 5.1.7 及其证明
（见 (5.1.13.1)）：
记
$I_{\vec\psi}
=I_{\psi^{(0)},\psi^{(1)}}$，
则
\begin{equation}
\mathcal A_{\ell_0}
\bigl(
[I_{\vec\psi}^{-1}\ell_q,
 I_{\vec\psi}^{-1}w']
\bigr)
=
\mathcal A_{\ell_0^{\vec\psi}}
([\ell_q,w']).
\tag{5.3.48}\label{eq:5-3-48}
\end{equation}
由 (5.3.18.4) 与 (5.3.48)，
\begin{equation}
\begin{aligned}
&
\mathcal A_{\ell_0^{\vec\psi}}
([\ell_q,w'])
-\mathcal A_{\ell_0}([\ell_p,w])
\\
&=
\mathcal A_{\ell_0}
\bigl([I_{\vec\psi}^{-1}\ell_q,
I_{\vec\psi}^{-1}w']\bigr)
-\mathcal A_{\ell_0}([\ell_p,w])
\\
&=
\mathcal A_{\ell_0}
\bigl([I_{\vec\psi}^{-1}\ell_q,
\widetilde u]\bigr)
=\int\widetilde u^{*}\omega.
\end{aligned}
\tag{5.3.49}\label{eq:5-3-49}
\end{equation}
此处
\[
\widetilde u(\rho,t)
=
\bigl(
\psi^{(1)}_{\rho t}
\circ
\psi^{(0)}_{\rho(1-t)}
\bigr)^{-1}
\bigl(\bar u(\rho,t)\bigr).
\]
定义微分同胚族
\begin{equation}
\psi_{\rho,t}
=
\bigl(
\psi^{(1)}_{\rho t}
\circ
\psi^{(0)}_{\rho(1-t)}
\bigr)^{-1}.
\tag{5.3.50}\label{eq:5-3-50}
\end{equation}
于是
$\widetilde u(\rho,t)
=\psi_{\rho,t}(\bar u(\rho,t))$。\translatornote{原书实体页 332
在 (5.3.50) 后写
$\psi_{\rho,t}
(\bar u(\chi^{-1}(\rho),t))$，
但 $\bar u$ 的首参数属于 $(0,1)$；
在 (5.3.52) 第二行又把
$\bar u$ 写成 $u$。
这里均依 $\widetilde u=\psi\circ\bar u$
的定义排印，见疑误 A-0306。}
现在计算
\begin{equation}
\begin{aligned}
\int\widetilde u^*\omega
&=\int_0^1\!\!\int_0^1
\omega\!\left(
\frac{\partial\widetilde u}{\partial\rho},
\frac{\partial\widetilde u}{\partial t}
\right)d\rho\,dt
\\
&=\int_0^1\!\!\int_0^1
\omega\!\left(
\frac{\partial}{\partial\rho}
\bigl(\psi_{\rho,t}(\bar u(\rho,t))\bigr),
\frac{\partial}{\partial t}
\bigl(\psi_{\rho,t}(\bar u(\rho,t))\bigr)
\right)d\rho\,dt.
\end{aligned}
\tag{5.3.51}\label{eq:5-3-51}
\end{equation}
又有
\begin{equation}
\begin{aligned}
\frac{\partial\widetilde u}{\partial\rho}
&=
\frac{\partial}{\partial\rho}
\bigl(\psi_{\rho,t}(\bar u(\rho,t))\bigr)
=X_{K_{\rho,t}}(\widetilde u)
+d\psi_{\rho,t}
\left(\frac{\partial\bar u}{\partial\rho}
(\rho,t)\right),
\\
\frac{\partial\widetilde u}{\partial t}
&=
\frac{\partial}{\partial t}
\bigl(\psi_{\rho,t}(\bar u(\rho,t))\bigr)
=X_{F_{\rho,t}}(\widetilde u)
+d\psi_{\rho,t}
\left(\frac{\partial\bar u}{\partial t}
(\rho,t)\right),
\end{aligned}
\tag{5.3.52}\label{eq:5-3-52}
\end{equation}
其中 $K_{\rho,t}$、
$F_{\rho,t}$ 分别是生成
$\frac{\partial\psi_{\rho,t}}{\partial\rho}
\circ\psi_{\rho,t}^{-1}$、
$\frac{\partial\psi_{\rho,t}}{\partial t}
\circ\psi_{\rho,t}^{-1}$
的哈密顿函数。

于是，
\begin{equation}
\begin{aligned}
\int\widetilde u^*\omega
&=\int_0^1\!\!\int_0^1
\omega\Bigl(
X_{K_{\rho,t}}(\widetilde u)
+d\psi_{\rho,t}
\bigl(\partial_\rho\bar u(\rho,t)\bigr),
\\[-1mm]
&\hspace{40mm}
X_{F_{\rho,t}}(\widetilde u)
+d\psi_{\rho,t}
\bigl(\partial_t\bar u(\rho,t)\bigr)
\Bigr)d\rho\,dt
\\
&=\int_{[0,1]^2}
\omega\!\left(
\partial_\rho\widetilde u,
X_{F_{\rho,t}}(\widetilde u)
\right)
+\int_{[0,1]^2}
\omega\!\left(
X_{K_{\rho,t}}(\widetilde u),
\partial_t\widetilde u
\right)
\\
&\quad-
\int_{[0,1]^2}
\omega\!\left(
X_{K_{\rho,t}}(\widetilde u),
X_{F_{\rho,t}}(\widetilde u)
\right)
+\int\bar u^*\omega.
\end{aligned}
\tag{5.3.53}\label{eq:5-3-53}
\end{equation}
这里用到了
(5.3.48)、(5.3.51)、(5.3.52)，
恒等式
\[
(a+b)(c+d)
=(a+b)d+b(c+d)-bd+ac
\]
以及 $\psi_{\rho,t}$ 的辛性。
还应注意
\[
\int_{[0,1]^2}\bar u^*\omega
=\int_{\mathbb R\times[0,1]}u^*\omega.
\]
有限能量条件保证此积分收敛；
若 $u$ 满足柯西--黎曼方程
(5.3.18.1)，这一项非负。

现在计算 (5.3.53)
右端的前三项。
由哈密顿向量场的定义，
\begin{equation}
\begin{aligned}
&\int_0^1\!\!\int_0^1
\omega\!\left(
\partial_\rho\widetilde u,
X_{F_{\rho,t}}(\widetilde u)
\right)d\rho\,dt
\\
&\quad=
-\int_0^1\!\!\int_0^1
dF_{\rho,t}
\left(\partial_\rho\widetilde u\right)
d\rho\,dt
\\
&\quad=
-\int_0^1
\bigl(
F_{1,t}(\widetilde u(1,t))
-F_{0,t}(\widetilde u(0,t))
\bigr)dt
+\int_0^1\!\!\int_0^1
\frac{\partial F_{\rho,t}}{\partial\rho}
(\widetilde u)\,d\rho\,dt.
\end{aligned}
\tag{5.3.54}\label{eq:5-3-54}
\end{equation}
同理，
\begin{equation}
\begin{aligned}
&\int_0^1\!\!\int_0^1
\omega\!\left(
X_{K_{\rho,t}}(\widetilde u),
\partial_t\widetilde u
\right)d\rho\,dt
\\
&\quad=
\int_0^1
\bigl(
K_{\rho,1}(\widetilde u(\rho,1))
-K_{\rho,0}(\widetilde u(\rho,0))
\bigr)d\rho
\\
&\qquad-
\int_0^1\!\!\int_0^1
\frac{\partial K_{\rho,t}}{\partial t}
(\widetilde u)\,d\rho\,dt.
\end{aligned}
\tag{5.3.55}\label{eq:5-3-55}
\end{equation}
最后，
\begin{equation}
\begin{aligned}
&-\int_0^1\!\!\int_0^1
\omega\!\left(
X_{K_{\rho,t}}(\widetilde u),
X_{F_{\rho,t}}(\widetilde u)
\right)d\rho\,dt
\\
&\qquad=
\int_0^1\!\!\int_0^1
\{F_{\rho,t},K_{\rho,t}\}
\,d\rho\,dt.
\end{aligned}
\tag{5.3.56}\label{eq:5-3-56}
\end{equation}
将 (5.3.54)--(5.3.56)
相加，再用恒等式 (5.3.43)，
便得到
\begin{equation}
\begin{aligned}
\int\widetilde u^*\omega
&=
-\int_0^1
\bigl(
F_{1,t}(\widetilde u(1,t))
-F_{0,t}(\widetilde u(0,t))
\bigr)dt
\\
&\quad+
\int_0^1
\bigl(
K_{\rho,1}(\widetilde u(\rho,1))
-K_{\rho,0}(\widetilde u(\rho,0))
\bigr)d\rho
+\int u^*\omega.
\end{aligned}
\tag{5.3.57}\label{eq:5-3-57}
\end{equation}
还需把 (5.3.57) 的前两项
与 $\|\psi^{(0)}\|$、
$\|\psi^{(1)}\|$
两段霍费尔长度联系起来。
设 $H^{(0)},H^{(1)}$
分别为同痕
$\psi^{(0)},\psi^{(1)}$
的归一化生成哈密顿函数。

首先，
$F_{1,t}$ 生成同痕
$\psi_{1,t}
=(\psi_t^{(1)})^{-1}$，
故
\begin{equation}
F_{1,t}(x)
=-H_t^{(1)}
\bigl(\psi_t^{(1)}(x)\bigr).
\tag{5.3.58}\label{eq:5-3-58}
\end{equation}
类似地，
\[
\psi_{0,t}
=(\psi_t^{(0)})^{-1},
\qquad
\psi_{\rho,0}\equiv\mathrm{id},
\qquad
\psi_{\rho,1}
=\bigl(
\psi_\rho^{(1)}
\circ\psi_{1-\rho}^{(0)}
\bigr)^{-1}.
\]
按哈密顿函数的复合与逆映射公式，
\begin{equation}
\begin{aligned}
F_{0,t}(x)
&=-H_t^{(0)}
\bigl(\psi_t^{(0)}(x)\bigr),\\
K_{\rho,0}(x)&=0,\\
K_{\rho,1}(x)
&=
H_{1-\rho}^{(0)}
\bigl(\psi_{1-\rho}^{(0)}(x)\bigr)
-H_\rho^{(1)}
\bigl(
\psi_\rho^{(1)}
(\psi_{1-\rho}^{(0)}(x))
\bigr).
\end{aligned}
\tag{5.3.59}\label{eq:5-3-59}
\end{equation}
把 (5.3.58)、(5.3.59)
代入 (5.3.57)，可得
\begin{equation}
\begin{aligned}
&-\int_0^1
\bigl(
F_{1,t}(\widetilde u(1,t))
-F_{0,t}(\widetilde u(0,t))
\bigr)dt
\\
&\qquad\ge
\int_0^1
\left(
\min_x H_t^{(1)}
-\max_x H^{(0)}
\right)dt,
\end{aligned}
\tag{5.3.60}\label{eq:5-3-60}
\end{equation}
且
\begin{equation}
\begin{aligned}
&\int_0^1
\bigl(
K_{\rho,1}(\widetilde u(\rho,1))
-K_{\rho,0}(\widetilde u(\rho,0))
\bigr)d\rho
\\
&\quad\ge
\int_0^1
\left(
\min_x H_{1-\rho}^{(0)}
-\max_x H_\rho^{(1)}
\right)d\rho
\\
&\quad=
\int_0^1
\left(
\min_x H_\rho^{(0)}
-\max_x H_\rho^{(1)}
\right)d\rho.
\end{aligned}
\tag{5.3.61}\label{eq:5-3-61}
\end{equation}
现在再假设 $u$ 满足
柯西--黎曼方程 (5.3.18.1)。
将 (5.3.60)、(5.3.61)
及 $\int u^*\omega\ge0$
代入 (5.3.57)，得到
\begin{equation}
\int\widetilde u^*\omega
\ge-
\bigl(
\|\psi^{(0)}\|
+\|\psi^{(1)}\|
\bigr).
\tag{5.3.62}\label{eq:5-3-62}
\end{equation}
证明完成。
\end{proof}

\begin{remark}{5.3.63}\label{rem:5-3-63}
若固定端点哈密顿微分同胚
$\psi_1^{(0)}$、
$\psi_1^{(1)}$，
可以遍取具有相同端点的同痕，
取下确界
\[
\begin{aligned}
\|\psi_1^{(1)}\|
&=
\inf_{H^{(1)}\mapsto\psi^{(1)}}
\|H^{(1)}\|,
\\
\|\psi_1^{(0)}\|
&=
\inf_{H^{(0)}\mapsto\psi^{(0)}}
\|H^{(0)}\|.
\end{aligned}
\]
它们分别是
$\psi_1^{(1)}$、
$\psi_1^{(0)}$ 的霍费尔范数。
如此可使常数
$\mathfrak E
=\mathfrak E(\psi^{(0)},\psi^{(1)})$
任意接近
$\|\psi_1^{(0)}\|
+\|\psi_1^{(1)}\|$。
然而，模空间 (5.3.19)
依赖同痕
$\psi^{(0)},\psi^{(1)}$
的选择，
所以这一能量估计
不能只用最终的
哈密顿微分同胚来表述。
如塞德尔在 \textbf{[Sei03I]}
所做的，
可以用条带
$\mathbb R\times[0,1]$
上满足拉格朗日边界条件的
哈密顿纤维丛全纯截面
重述弗洛尔理论；
在此框架下，
这类能量估计可以
更为内在地表述。
\end{remark}
